%========================================================================================
\section{Environnement et technologies}
\label{sec:environnement}
%========================================================================================

\subsection{Versions retenues}

Le tableau~\ref{tab:versions} donne les versions fixées dans \code{requirements.txt} et \code{requirements-dev.txt}. Chaque dépendance a une version exacte. Deux installations donnent donc le même environnement.

\begin{table}[H]
  \centering
  \small
  \caption{Versions des outils et des bibliothèques.}
  \label{tab:versions}
  \begin{tabularx}{\linewidth}{@{}l l X@{}}
    \toprule
    \textbf{Composant} & \textbf{Version} & \textbf{Rôle} \\
    \midrule
    Python & 3.14 & Langage \\
    PostgreSQL & 18 & Base de données relationnelle~\cite{postgresql} \\
    Django & 5.2.18 (LTS) & Cadriciel web, ORM, administration~\cite{django} \\
    Django REST Framework & 3.18.3 & API REST~\cite{drf} \\
    SimpleJWT & 5.5.1 & Jetons JWT et liste noire des jetons révoqués~\cite{simplejwt} \\
    drf-spectacular & 0.30.0 & Schéma OpenAPI 3, Swagger UI, ReDoc~\cite{spectacular} \\
    Graphene-Django & 3.2.3 & API GraphQL~\cite{graphene} \\
    django-filter & 26.2 & Filtres des listes \\
    django-cors-headers & 4.9.0 & Politique CORS \\
    django-environ & 0.14.0 & Configuration par variables d'environnement \\
    psycopg & 3.3.6 & Pilote PostgreSQL \\
    requests & 2.34.2 & Appels HTTP vers les services externes \\
    Gunicorn, WhiteNoise & 26.2.0, 6.12.0 & Serveur WSGI et fichiers statiques en production \\
    \midrule
    pytest, pytest-django & 9.1.1, 4.14.0 & Tests \\
    pytest-cov & 7.1.0 & Couverture de code \\
    responses & 0.26.3 & Simulation des appels HTTP externes \\
    Ruff & 0.17.1 & Analyse statique et formatage \\
    \bottomrule
  \end{tabularx}
\end{table}

Nous avons développé sous Windows 11. Les scripts d'exploitation sont donc en PowerShell, et nous avons géré la base avec pgAdmin~4.

\subsection{Pourquoi Django 5.2 LTS}

Nous voulions la dernière version stable de Django qui marche avec \emph{toutes} nos dépendances. Django 6.1 était déjà sorti, mais Graphene-Django 3.2.3 et SimpleJWT 5.5.1 ne disent pas encore supporter Django~6. Django 5.2 est une version LTS, maintenue jusqu'en avril 2028, et elle supporte officiellement Python 3.14. Nous avons préféré un ensemble stable à la version la plus récente du cadriciel.

\subsection{Configuration}

Les paramètres sont répartis dans quatre modules (tableau~\ref{tab:settings}). Les valeurs sensibles ne sont jamais écrites dans le code. \code{django-environ} les lit dans l'environnement, sinon dans un fichier \code{.env} que Git ignore. Toutes les variables sont expliquées dans \code{.env.example}.

\begin{table}[H]
  \centering
  \small
  \caption{Modules de paramètres.}
  \label{tab:settings}
  \begin{tabularx}{\linewidth}{@{}l X@{}}
    \toprule
    \textbf{Module} & \textbf{Usage} \\
    \midrule
    \code{config.settings.base} & Paramètres communs (applications, DRF, JWT, GraphQL, journaux) \\
    \code{config.settings.dev} & Développement, avec \code{DEBUG}, GraphiQL et l'introspection activés \\
    \code{config.settings.test} & Tests, avec un hachage rapide des mots de passe, un cache en mémoire et les journaux sur la console \\
    \code{config.settings.prod} & Production, avec HTTPS, HSTS, cookies sécurisés, JSON seulement et WhiteNoise \\
    \bottomrule
  \end{tabularx}
\end{table}
